\vspace{-.5em}
Understanding temporal dynamics in clinical narratives is essential for modeling patient trajectories, yet large-scale temporally annotated resources remain limited. We present PMOA–TTS, the first openly available dataset of 124,699 PubMed Open Access (PMOA) case reports, each converted into structured (event, time) timelines via a scalable LLM-based pipeline. Our approach combines heuristic filtering with Llama 3.3 to identify single-patient case reports, followed by prompt-driven extraction using Llama 3.3 and DeepSeek R1, resulting in over 5.6 million timestamped clinical events. To assess timeline quality, we evaluate against a clinician-curated reference set using three metrics: (i) event-level matching (80\% match at a cosine similarity threshold of 0.1), (ii) temporal concordance (c-index > 0.90), and (iii) Area Under the Log-Time CDF (AULTC) for timestamp alignment. Corpus-level analysis shows wide diagnostic and demographic coverage. In a downstream survival prediction task, embeddings from extracted timelines achieve time-dependent concordance indices up to \(0.82 \pm 0.01\), demonstrating the predictive value of temporally structured narratives. PMOA–TTS provides a scalable foundation for timeline extraction, temporal reasoning, and longitudinal modeling in biomedical NLP. The dataset is available at: \url{https://huggingface.co/datasets/snoroozi/pmoa-tts}.\vspace{-.5em}


% Understanding the temporal dynamics of clinical events in narrative text is critical for modeling patient trajectories, yet large-scale resources with temporal annotations remain scarce. We introduce PMOA–TTS, the first openly available dataset of 124,699 PubMed Open Access (PMOA) case reports, each transformed into structured timelines of (event, time) tuples using a scalable pipeline built on large language models (LLMs). Our method integrates heuristic filtering with Llama 3.3 to identify single-patient case reports, followed by prompt-driven extraction using Llama 3.3 and DeepSeek, yielding over 5.6 million timestamped clinical events. To evaluate timeline quality, we benchmark against a clinician-curated reference set using three complementary metrics: (i) event-level matching (achieving 80\% accuracy at a cosine similarity threshold of 0.1), (ii) temporal concordance (c-index > 0.90), and (iii) a novel Area Under the Log-Time CDF (AULTC) metric for timestamp alignment. Corpus-level analysis reveals broad diagnostic coverage and balanced demographic distributions. In a downstream survival prediction task, language model embeddings derived from extracted timelines achieve time-dependent concordance indices up to 
% \(0.82 \pm 0.01\), highlighting the utility of temporally structured clinical narratives for predictive modeling. PMOA–TTS establishes a scalable foundation for timeline extraction, temporal reasoning, and longitudinal modeling in biomedical NLP. \textcolor{red}{The code and data are available at <url> }